\section{共同犯罪 - 共同犯罪的分类、正犯与共犯的关系、间接正犯}
\begin{enumerate}
  \item \zd{人多不要怕，一个一个往里加。永远先分析实行犯。4 大角色 2 个阶层同时分析同时兼顾}。
  \item \al{甲教唆乙，乙帮助丙，丙实行抢劫了丁}。
  \item 分类。
        \begin{enumerate}
          \item \zd{从分工}。从侵害法益的方式上分。
                \begin{enumerate}
                  \item 实行犯/正犯。
                        \begin{enumerate}
                          \item 从数量上。单独正犯、\zd{共同正犯}（比较难）。
                          \item 从方式。直接正犯（自己上）、\zd{间接正犯}（幕后）。\al{甲让乙去盗窃，乙是甲盗窃的工具}。
                        \end{enumerate}
                  \item 共犯。
                        \begin{enumerate}
                          \item \zd{狭义共犯 = 教唆 + 帮助}。
                          \item 广义共犯 = 教唆 + 帮助 + 实行。
                        \end{enumerate}
                \end{enumerate}
          \item 从作用。主要是涉及共同犯罪的处罚问题。考的很少。
                \begin{enumerate}
                  \item 主犯。起主要作用。
                  \item 从犯。起次要作用。\al{所有的帮助犯都是从犯 \dui}。
                  \item 胁从犯。被迫参加犯罪。人身完全被强制、意志自由完全被剥夺 \yjt 紧急避险。
                        \al{甲掌握了乙女的裸照，威胁乙去教唆丙法官徇私枉法，乙被迫答应}。
                \end{enumerate}
          \item 从结构。
                \begin{enumerate}
                  \item 任意的共同犯罪，人数 \dydy1。
                  \item 必要的共同犯罪，人数 \dydy2。对合/对向犯。缺少一方，另一方无法单独构成犯罪。
                        \begin{enumerate}
                          \item 罪名、法定刑相同。\al{重婚。乙明知甲以结婚，仍与其领取结婚证}。
                          \item 罪名、法定刑不同。\al{行贿、受贿。甲送给国家工作人员乙 10W，托其将儿子录为公务员}。
                        \end{enumerate}
                  \item 片面的对向犯。只处罚一方。
                        \begin{enumerate}
                          \item 另一方的存在只是一方构成犯罪的条件。
                          \item \al{贩卖淫秽物品牟利、破坏军婚}。
                        \end{enumerate}
                \end{enumerate}
          \item 犯罪集团。
                \begin{enumerate}
                  \item 3 人以上为实施一次或多次犯罪，经事前通谋而建立的犯罪组织。是共同犯罪的一种特殊形式。
                  \item 首要分子都是主犯。按集团所犯的全部罪行处罚。
                \end{enumerate}
        \end{enumerate}
  \item \kd{共同正犯}。\zd{一起（纽带是意思联络，不一起 \yjt 片面共犯）实施实行行为}。喜欢考主观阶层步调不一致。
        \begin{enumerate}
          \item 责任年龄、责任能力。一个有，一个没有。
                \begin{enumerate}
                  \item \al{甲乙共同伤害丙，乙是精神病患者但甲不知}。
                  \item \al{甲教唆乙伤害丙，乙是精神病患者但甲不知}。
                \end{enumerate}
          \item \zd{对各自的主观（故意、过失）是否做要求}。
                \begin{enumerate}
                  \item 完全犯罪共同说。
                  \item 部分犯罪共同说（实务采取）。
                  \item \zd{行为共同说}。缺陷。不符合我国刑法规定，违反罪刑法定原则。\al{二人以上共同过失犯罪，不以共同犯罪论处}。
                \end{enumerate}
                \begin{kk}{无法查明的 4 种情形}
                  \begin{enumerate}
                    \item 故意的共同正犯。\al{甲乙一起枪杀丙 \yjt 共同犯罪，故意杀人既遂}。
                    \item 过失的共同正犯。\al{甲乙一起开枪打野猪，打死了丙}。
                          \begin{enumerate}
                            \item 部分犯罪共同说。过失 + 过失不构成共同犯罪 \yjt 必须都是故意 \yjt 单独定罪 \yjt 存疑时有利于被告 \yjt 都定未遂。
                            \item 行为共同说。\zd{部分实行、全部负责} \yjt 无法查明，无需查明 \yjt 过失的共同犯罪 \yjt 都定过失致人死亡。
                          \end{enumerate}
                    \item 故意的同时犯。\al{甲乙无意思联络，打死了丙 \yjt 不是共同犯罪，都定故意杀人未遂}。
                    \item 过失的同时犯。\al{甲乙无意思联络、无过失，打死了丙 \yjt 不是共同犯罪 \yjt 存疑时有利于被告 \yjt 都无罪}。
                  \end{enumerate}
                \end{kk}
        \end{enumerate}
  \item \kd{间接正犯} = \zd{引起 + 支配。A + B}。是否形成支配力，\zd{看实质条件}）。\zd{支配力来源}。
        \begin{enumerate}
          \item 强制手段。\al{赐死}。
          \item 欺骗手段。\al{欺骗不可以形成支配力 \cuo}。
                \begin{enumerate}
                  \item 欺骗他人，\zd{利用}他人有\zd{过失}的行为。
                        \al{医生骗护士打毒针。护士过失致人死亡 \hz 医疗事故，医生故意杀人间接正犯}。
                  \item 欺骗他人，\zd{利用}他人有犯罪\zd{故意}的行为。
                        \al{甲想杀在屏风后的乙，骗丙砸屏风。
                          丙故意毁坏财物、过失致人死亡（如果存在过失），想象竞合择一重。
                          甲故意毁坏财物教唆，故意杀人间接正犯}。
                  \item \zd{直接欺骗被害人实施自损行为}。
                        \al{甲挖坑杀乙，丙利用这个坑，骗丁走入，结果杀死了丙}。
                        \al{教唆精神病患者自杀、诱骗相约自杀 \yjt 故意杀人间接正犯}。
                \end{enumerate}
          \item \zd{法定身份}。拟制的支配力。
                \al{狱警甲指使囚犯乙虐待囚犯丙。乙故意伤害、虐待被监管人员帮助犯；甲故意伤害教唆、虐待被监管人员间接正犯}。
                破除早期间接正犯是一个单独犯罪的现象，和被利用人是不可能构成共同犯罪的。
          \item 命题陷阱。\zd{把过失类罪名安到间接正犯头上}。\al{甲应按交通肇事罪的间接正犯论处 \cuo}。
        \end{enumerate}
  \item \kd{正犯与共犯关系}。观点展示。
        \begin{enumerate}
          \item \zd{共犯从属性原理}。
                \begin{enumerate}
                  \item 实行者实行行为对法益的侵害最大、最直接，教唆者、帮助者是间接侵害，\zd{结果取决于实行者}。
                  \item 教唆者、帮助者要构成犯罪，必须\zd{以实行者构成犯罪为前提}。
                  \item \zd{犯罪形态}也从属实行者。
                \end{enumerate}
          \item 共犯独立性（少数说）。缺点。
                \begin{enumerate}
                  \item 独立看待教唆、帮助行为。
                  \item 站在主观主义立场，过于看重帮助者、教唆者的犯罪故意。
                \end{enumerate}
        \end{enumerate}
\end{enumerate}
